\documentclass[10pt]{article}
\usepackage[margin=12mm]{geometry}
\usepackage{graphicx}
\usepackage{caption}
\pagestyle{empty}
\captionsetup{font=small,labelfont=bf}

\begin{document}
\begin{figure}[p]
\centering
\includegraphics[width=\textwidth]{fig2_ce_terminal_pareto_front.pdf}
\caption{\textbf{The terminal layer of the \emph{C.\@ elegans} lineage is
nearly Pareto optimal.} Terminal-cell identities defined by their final physical
positions and molecular states measured from protein expression were held fixed while
terminal cells were reassigned to candidate parents by minimum-cost bipartite
matching across varying weightings of combined travel and cell-state distance. The
resulting curve traces the Pareto front. Each objective is reported in
standard deviations of its first-cousin-shuffle null and translated so the
natural lineage (black cross) is at the origin. As lineage restrictions are
relaxed, the null distributions move progressively farther into the
upper-right: first-cousin (gold), second-cousin (ochre), third-cousin
(brown), and full-random shuffling (plum; inset as it lies further beyond the
main-axis ranges). Pareto assignments are colored by edge retention, the
fraction of natural terminal parent--child edges preserved; lighter blue
denotes greater retention hence more preserved natural lineage structure.
The outlined circle identifies the Pareto assignment with maximum edge retention.}
\label{fig:ce_terminal_pareto_main}
\end{figure}
\end{document}
